\subsection*{李群的例子}

\begin{enumerate}
\def\labelenumi{(\arabic{enumi})}
\tightlist
\item
  令 \(\mathbf{E}\) 为 \(\mathbf{K}\) 上的完备可赋范空间。映射 \((x, y) \mapsto x - y\) 从 \(\mathbf{E} \times \mathbf{E}\) 到 \(\mathbf{E}\)，是连续线性的，因而是解析的。
\end{enumerate}

因此，E 配以其加法群结构和解析流形结构后是李群。

特别地，K 是李群。

\begin{enumerate}
\def\labelenumi{(\arabic{enumi})}
\setcounter{enumi}{1}
\tightlist
\item
  令 A 为 K 上的完备可赋范含单位结合代数。从 A × A 到 A 的乘法 \((x,y)\mapsto xy\) 是双线性且连续的，因而是解析的。命题 3 表明，A 的可逆元群 A* 在 A 中是开集（这也可由 General Topology, Chapter IX, §3, Proposition 13 得到），并且 A* 是李群。
\end{enumerate}

例如，设 E 为 K 上的完备可赋范空间，并令 A = \(\mathcal{L}(\mathrm{E})\)（General Topology, Chapter IX, § 3, Proposition 5）。于是 A* 是 E 的自同构群 GL(E)。因此，这个群自然地带有 K 上的李群结构。更具体地说，GL(n, K) 配以由 \(\mathbf{M}_{n}(\mathbf{K})\) 诱导的流形结构后是李群。对于 \(n = 1\)，可见乘法群 K* 配以由 K 诱导的流形结构后是李群。

\begin{enumerate}
\def\labelenumi{(\arabic{enumi})}
\setcounter{enumi}{2}
\tightlist
\item
  令 G 为 K 上的李群。令 \(K' = R\) 或 C 或一个非离散完备超度量域，并令 \(\sigma\) 为赋值域 \(K'\) 到 K 的一个赋值子域上的同构。于是 G 配以通过限制标量得到的 K'-流形结构后，是 \(K'\) 上的李群；称其为由李群 G 通过限制标量导出的李群（从 K 到 \(K'\) 通过 \(\sigma\)）。例如，每个复李群都自然地带有实李群结构。此外，每个复李群 G 都关联着一个称为 G 的共轭的复李群，它由 C 的自同构 \(z \mapsto \bar{z}\) 从 G 导出。
\end{enumerate}

